\Article{14}
We still have to consider in a little more detail a circumstance that arises in our general solution. In Article 5 we showed that there are always two solutions. In one, \(P+iQ\) must be a function of \(p+iq\) and \(P-iQ\) a function of \(p-iq\). In the other, \(P+iQ\) must be a function of \(p-iq\) and \(P-iQ\) a function of \(p+iq\). We will now show that in one solution the parts of the representation always have a similar position to the corresponding parts of the original, while in the other their position is reversed. We will also give a criterion for distinguishing the two in advance.

First, we note that we can speak of perfect or reversed similarity only if we distinguish two sides on each surface and regard one as upper and the other as lower. Since this choice is in itself somewhat arbitrary, the two solutions do not differ essentially. A reversed similarity becomes a perfect one as soon as we take the side of one surface previously regarded as upper to be lower. This distinction could therefore play no part in our solution, since the surfaces were defined solely by the coordinates of their points. To take the distinction into account, we must first specify the surfaces in another way that incorporates it. For this purpose let the nature of the first surface be determined by the equation \(\psi=0\), where \(\psi\) is a given single-valued function of \(x,y,z\). Thus at every point of the surface the
\SourceLeaf{217}
value of \(\psi\) vanishes, and at every point of space outside the surface it is nonzero. On crossing the surface, at least generally speaking, the value of \(\psi\) passes from positive to negative, and on crossing in the opposite direction from negative to positive. In other words, the value of \(\psi\) is positive on one side and negative on the other. We take the former to be the upper side and the latter the lower side. We proceed in exactly the same way for the second surface, whose nature is determined by the equation \(\Psi=0\), where \(\Psi\) is a given single-valued function of the coordinates \(X,Y,Z\). Let differentiation further give
\[
\begin{aligned}
\mathrm{d}\psi&=e\,\mathrm{d}x+g\,\mathrm{d}y+h\,\mathrm{d}z\\
\mathrm{d}\Psi&=E\,\mathrm{d}X+G\,\mathrm{d}Y+H\,\mathrm{d}Z
\end{aligned}
\]
\noindent where \(e,g,h\) are functions of \(x,y,z\) and \(E,G,H\) are functions of \(X,Y,Z\).

The considerations needed to reach our goal are somewhat unusual, though they present no difficulty in themselves, so we will try to explain them as clearly as possible. Between the two corresponding representations on the surfaces with equations \(\psi=0\) and \(\Psi=0\), let us take six intermediate representations in the plane. We then have eight different representations to consider.
\begin{center}
\begingroup\fontsize{8}{10}\selectfont\setlength{\tabcolsep}{2pt}
\begin{tabular}{@{}lll@{}}
\multicolumn{3}{r}{\parbox{1.55in}{\raggedright taking as corresponding the points whose coordinates are, respectively, \(=\)}}\\[2pt]
1\textsuperscript{s}&the original on the surface with equation \(\psi=0\) . .&\(x,\ y,\ z\)\\
2\textsuperscript{s}&representation in the plane . . . . . . . . . . . . . . . . . . . . .&\(x,\ y,\ 0\)\\
3\textsuperscript{s}&\quad „\quad „\quad „\quad „\quad . . . . . . . . . . . . . . . . . . . . .&\(t,\ u,\ 0\)\\
4\textsuperscript{s}&\quad „\quad „\quad „\quad „\quad . . . . . . . . . . . . . . . . . . . . .&\(p,\ q,\ 0\)\\
5\textsuperscript{s}&\quad „\quad „\quad „\quad „\quad . . . . . . . . . . . . . . . . . . . . .&\(P,\ Q,\ 0\)\\
6\textsuperscript{s}&\quad „\quad „\quad „\quad „\quad . . . . . . . . . . . . . . . . . . . . .&\(T,\ U,\ 0\)\\
7\textsuperscript{s}&\quad „\quad „\quad „\quad „\quad . . . . . . . . . . . . . . . . . . . . .&\(X,\ Y,\ 0\)\\
8\textsuperscript{s}&representation on the surface with equation \(\Psi=0\) . .&\(X,\ Y,\ Z\)
\end{tabular}\endgroup
\end{center}
\par\indent We will now compare these representations solely with respect to the relative \textit{position} of the infinitely small linear elements, %
\SourceLeaf{218}%
leaving their relative magnitudes entirely aside. We regard two representations as similarly positioned if, for two linear elements issuing from one point, the element on the right in one representation corresponds to the element on the right in the other. Otherwise we call their positions reversed. For the plane in Nos. 2---7 we always regard the side on which the third coordinate takes positive values as the upper side. For the first and last surfaces, however, the distinction between upper and lower depends solely on whether \(\psi\) and \(\Psi\) are positive or negative, as specified above.

First, it is clear that at any point of the first surface where, with \(x\) and \(y\) fixed, a positive increment in \(z\) takes us to its upper side, representation 2 has a position similar to that of representation 1. This evidently holds wherever \(h\) is positive. The reverse occurs when \(h\) is negative, and the representations have reversed positions.

In the same way, representations 7 and 8 are similarly positioned or reversed, according as \(H\) is positive or negative.

To compare representations 2 and 3, let \(\mathrm{d}s\) in the former be the length of an infinitely small line from the point with coordinates \(x,y\) to another with coordinates \(x+\mathrm{d}x,y+\mathrm{d}y\), and let \(l\) be its inclination to the abscissa axis, increasing in the direction in which we turn from the \(x\)-axis to the \(y\)-axis. Thus
\[
\mathrm{d}x=\mathrm{d}s\mathord{.}\cos l,\quad \mathrm{d}y=\mathrm{d}s\mathord{.}\sin l
\]
\noindent In representation 3 let \(\mathrm{d}\sigma\) be the length of the line corresponding to \(\mathrm{d}s\), and let its inclination to the abscissa axis, in the same sense as before, be \(\lambda\), so that
\[
\mathrm{d}t=\mathrm{d}\sigma\mathord{.}\cos\lambda,\quad \mathrm{d}u=\mathrm{d}\sigma\mathord{.}\sin\lambda
\]
\noindent Hence, in the notation of Article 4, we have
\[
\begin{aligned}
\mathrm{d}s\mathord{.}\cos l&=\mathrm{d}\sigma(a\cos\lambda+a'\sin\lambda)\\
\mathrm{d}s\mathord{.}\sin l&=\mathrm{d}\sigma(b\cos\lambda+b'\sin\lambda)
\end{aligned}
\]
\noindent and therefore
\[
\mathrm{tang}\,l=\frac{b\cos\lambda+b'\sin\lambda}{a\cos\lambda+a'\sin\lambda}
\]
